\documentclass{article}
\usepackage[utf8]{inputenc}
\usepackage{amsmath}
\usepackage{geometry}
\geometry{margin=2cm}
\begin{document}

\section*{Questão 96 — ENEM 2024 — Ciências da Natureza}

Os pesticidas naturais vêm sendo utilizados no controle de pragas e doenças agrícolas como substituintes de pesticidas sintéticos tradicionais, por serem menos nocivos ao ambiente, biodegradáveis e minimizarem custos e riscos relativos à lavoura. Por exemplo, os compostos 1 e 2 apresentados estão envolvidos nas respostas de defesa das plantas. Os grupos funcionais presentes nesses compostos são importantes para suas propriedades no controle de pragas.

\textbf{Pergunta:} Qual é a função orgânica correspondente ao grupo funcional comum presente nesses dois compostos?

\section*{Interpretação do problema e estratégia}

A questão apresenta estruturas químicas dos compostos naturais e solicita a identificação da função orgânica correspondente ao grupo funcional comum a ambos.

\begin{itemize}
  \item Identificar grupos funcionais em cada composto.
  \item Confrontar as estruturas para determinar o grupo comum.
  \item Selecionar a função orgânica correta entre as alternativas disponíveis.
\end{itemize}

\section*{Mini revisão de conteúdo}

\begin{tabular}{|c|p{6cm}|p{7cm}|}
\hline
Função Orgânica & Descrição & Observação na questão \\
\hline
Ácido carboxílico & Possui grupo \( \ce{-COOH} \) formado por carbonila \(\ce{C=O}\) e hidroxila \(\ce{OH}\). & Presente em ambos os compostos -- grupo funcional comum.
\\
\hline
Cetona & Grupo carbonila \(\ce{C=O}\) sem hidroxila, ligada a dois radicais. & Não é comum.
\\
\hline
Fenol & Hidroxila ligada diretamente ao anel aromático. & Só no composto 1.
\\
\hline
\end{tabular}

\par
\textbf{Considerações:} Diferenciar claramente ácido carboxílico de cetona pela presença do grupo OH, importante para propriedades químicas e biológicas.

\section*{Resolução}

\begin{enumerate}
  \item Analisando o Composto 1: grupo ácido carboxílico e fenol presentes.
  \item Composto 2 também possui o grupo ácido carboxílico.
  \item Portanto, o grupo funcional comum é o ácido carboxílico \( \ce{-COOH} \).
  \item A alternativa correta é a \textbf{A} (ácido carboxílico).
\end{enumerate}

\textbf{Pulo do gato:} Foque no grupo funcional recorrente em ambas as estruturas, que é o ácido carboxílico.

\section*{Armadilhas e erros comuns}

\begin{itemize}
  \item Confundir cetona com ácido carboxílico apenas por presença de carbonila.
  \item Achar que o grupo hidroxila do composto 1 é álcool comum e que aparece também no composto 2.
  \item Pensar que a ligação dupla C=C é a função comum.
  \item Considerar fenol como grupo comum, que só ocorre no composto 1.
\end{itemize}

\section*{Código da questão}

\texttt{CIENCIAS\_DA\_NATUREZA\_ORGANICA\_001}

\bigskip

\textbf{\large Resolução elaborada baseada no enunciado e na análise estrutural dos compostos químicos apresentados.}

\end{document}